\section{来源审计：把个人案件转化为组织工程课}

Joe Sullivan 的职业路径横跨联邦网络犯罪检察、科技公司 security leadership、漏洞披露生态与公益项目。课堂因此同时包含职业经验、个人案件叙述、政策判断和领导力建议。高质量讲义不能把这些内容混成一种“事实”：讲者对事件的回忆属于课堂陈述，检方的主张属于指控，陪审团与法院处理的是法律责任，而本文负责提炼可复用的工程和治理机制。

\begin{warningbox}{法律与时间边界}
本文是 security engineering 与 governance 讲义，不构成法律意见。课堂录制于 2025 年 1 月，当时 Sullivan 的上诉尚未裁决；第九巡回法院于 2025 年 3 月 13 日维持定罪，2025 年 11 月 12 日发布修订意见并拒绝全院复审，美国最高法院于 2026 年 6 月 29 日拒绝受理 certiorari petition。课堂中的上诉预期只能作为历史观点，不能写成当前程序状态。
\end{warningbox}

\begin{importantbox}{本讲的三个治理闭环}
\begin{enumerate}
  \item \textbf{External-report loop}：接收外部报告 $\rightarrow$ 验证与定级 $\rightarrow$ 修复 $\rightarrow$ 协调披露或奖励；
  \item \textbf{Incident loop}：检测 $\rightarrow$ 控制影响 $\rightarrow$ 保存证据 $\rightarrow$ 调查、恢复与复盘；
  \item \textbf{Disclosure loop}：技术事实 $\rightarrow$ 法务与 materiality 分析 $\rightarrow$ 高管/董事会决策 $\rightarrow$ 对监管者、客户或公众沟通。
\end{enumerate}
三个闭环可以由同一条线索触发，却不能由同一个人、同一张 ticket 或同一种激励机制代替。
\end{importantbox}

\subsection{Security program 不是一个“英雄 CSO”}

安全领导者经常以个人能力被评价：能否发现攻击、能否说服 CEO、能否在危机里快速拍板。然而规模化组织真正需要的是 \term{security program}（安全项目）：明确人员能力、重复流程、技术平台和治理权责，使正确行为不依赖某位负责人临场发挥。Sullivan 从检察官转向 eBay、Facebook、Uber 等技术公司的经历，最有价值的共同主题不是公司名，而是安全团队必须随组织复杂度一起升级。

\lecturefigure{01-security-program.png}{可持续的 security program 同时依赖 people、process、platform 与 governance。}{本地字幕 00:00--03:20；概念重绘。}

\begin{knowledgebox}{读图：四层缺一不可}
先看左上角的 people：人才、值班、沟通和决策能力决定组织能否理解事件。Process 把 intake、incident command、disclosure 与复盘变成可重复路径。Platform 提供 identity、telemetry、case management 和 evidence storage。Governance 则回答谁有 authority、谁接受 challenge、谁签署外部披露。只建设工具会得到“告警很多、决定很少”；只依赖专家会得到“专家离开、系统失忆”。
\end{knowledgebox}

\teachervoice{Sullivan 的职业故事把法律、工程和组织放在同一个视野里。课堂提示：security leader 的稀缺价值不是比所有工程师更懂每个漏洞，而是能把技术事实翻译成法律、业务和公共风险，同时确保每种判断由有权负责的人作出。}

\begin{knowledgebox}{术语消化：角色与责任}
\term{CSO/CISO} 负责安全能力、风险信息和升级；\term{general counsel}（总法律顾问，GC）负责组织法律分析与律师管理；\term{CEO} 对企业优先级与重大经营决定负责；\term{board/disclosure committee} 提供监督并处理重大披露。角色可以因公司规模不同而合并，但“谁提供事实、谁给法律意见、谁承担经营决定”必须显式区分。
\end{knowledgebox}

\subsection{本章小结}

本讲不以“谁说得更有道理”为主线，而以 operating system 为主线：外部报告、事件响应与披露治理分别需要 owner、evidence、decision right 和记录。个人案件只用于检验这些机制为何必须存在。

